% ============================================================================
% templates/exam/tarea/tarea.tex — plantilla: tarea / trabajo domiciliario
% ----------------------------------------------------------------------------
% Para: tareas, asignaciones, trabajos aplicativos y proyectos breves con
% condiciones de entrega y criterios de calificación explícitos.
% Uso: scripts/build.sh tarea.tex   (LuaLaTeX)
% ============================================================================
\documentclass{academic-exam}

\universidad{Universidad Nacional de San Cristóbal de Huamanga}
\facultad{Facultad de Ciencias Económicas, Administrativas y Contables}
\escuela{Escuela Profesional de Economía}
\curso{Finanzas}
\codigocurso{EC-361}
\docente{Mg. Nombre Apellido}
\periodo{Semestre 2026-I}
\tipoevaluacion{Tarea 05}
\fechaevaluacion{Publicada: 21 de julio de 2026}
\duracion{Entrega: 4 de agosto de 2026, 23:59}
\modalidad{Domiciliaria · individual}

\begin{document}
\membrete

\begin{instrucciones}[Condiciones de entrega]
\begin{itemize}
  \item Entregue un único archivo PDF por el aula virtual, nombrado
        \texttt{tarea05-apellido-nombre.pdf}. Adjunte además la hoja de
        cálculo (\texttt{.xlsx} u \texttt{.ods}) con sus cálculos.
  \item Las entregas fuera de plazo se califican sobre 15 puntos hasta
        24 horas después; luego no se reciben.
  \item El trabajo es individual: entregas con similitud sustancial se
        califican con cero para todos los involucrados.
\end{itemize}
\end{instrucciones}

\begin{ejercicio}[6][Valor del dinero en el tiempo]
Una municipalidad evalúa dos ofertas para financiar maquinaria por
S/~600\,000:
\begin{apartados}
  \item Banco A: préstamo a 5 años, TEA de 11{,}5\,\%, cuotas anuales
        constantes. Construya la tabla de amortización completa. \pts{3}
  \item Banco B: TEA de 10{,}8\,\% con comisión inicial de 2\,\% del
        principal. Calcule el costo efectivo anual de cada oferta y
        recomiende una, justificando. \pts{3}
\end{apartados}
\end{ejercicio}

\begin{ejercicio}[7][Evaluación de proyectos]
Un proyecto agroindustrial requiere una inversión de S/~850\,000 y genera
flujos netos anuales de S/~210\,000 durante 6 años, con un valor de
rescate de S/~120\,000.
\begin{apartados}
  \item Calcule el VAN con un costo de oportunidad de 12\,\% y la TIR del
        proyecto. \pts{3}
  \item Elabore el análisis de sensibilidad del VAN ante variaciones del
        flujo anual de $\pm 10\,\%$ y $\pm 20\,\%$; presente los
        resultados en una tabla y un gráfico. \pts{2}
  \item Calcule el periodo de recuperación descontado. \pts{2}
\end{apartados}
\end{ejercicio}

\begin{ejercicio}[7][Aplicación con datos reales]
Descargue de la SMV (o de la BVL) los estados financieros 2023--2025 de
una empresa peruana listada, del sector de su elección.
\begin{apartados}
  \item Calcule e interprete los ratios de liquidez, gestión, solvencia y
        rentabilidad de los tres años. \pts{4}
  \item Redacte un diagnóstico financiero de una carilla que integre los
        ratios con la coyuntura del sector. \pts{3}
\end{apartados}
\end{ejercicio}

\begin{observacion}[Recomendaciones]
Presente los cálculos intermedios, no solo los resultados. En la hoja de
cálculo, use una pestaña por ejercicio y rotule las celdas: una hoja
desordenada que impida verificar el procedimiento pierde el puntaje del
apartado correspondiente.
\end{observacion}

\begin{criterios}
\begin{tabularx}{\linewidth}{@{}X c@{}}
  \textbf{Criterio} & \textbf{Puntaje}\\
  \arrayrulecolor{grislinea}\hline\addlinespace[3pt]
  Exactitud de cálculos y procedimientos           & 12\\
  Interpretación y diagnóstico                     & 5\\
  Presentación, formato y reproducibilidad         & 3\\
\end{tabularx}
\end{criterios}

\findelexamen
\end{document}
